% Source: https://artofproblemsolving.com/community/c3h2491537_help_with_understanding
% Stable identity: aops-topic-2491537-q01
Number the cars $k=0,1,\dots,50$ in order of departure. Car $k$ leaves at time $k$ hours with speed $50+k$, so at time $100$ hours it has been driving for $100-k$ hours and is at distance
\[d_k=(50+k)(100-k)=5000+50k-k^2=5625-(k-25)^2\]
kilometres from $A$. This is largest exactly when $k=25$, so the leading car is car $25$, with speed $50+25=\boxed{\textbf{(C)}\ 75}$ kilometres per hour.
